\Needspace{6\baselineskip}\Article{353}
\textit{First example, for \(n=19\).} Since here \(n-1=3.3.2\), finding the roots \(\Omega\) is to be reduced to solving two cubic equations and one quadratic equation. This example will be easier to understand because most of the necessary operations are already contained in the preceding articles. Taking the number 2 for the primitive root \(g\) gives the following least residues of its powers (the exponents of the powers are written above the residues in the first row): {\small\setlength{\arraycolsep}{2pt}\[\begin{array}{rrrrrrrrrrrrrrrrrr}0.&1.&2.&3.&4.&5.&6.&7.&8.&9.&10.&11.&12.&13.&14.&15.&16.&17\\1.&2.&4.&8.&16.&13.&7.&14.&9.&18.&17.&15.&11.&3.&6.&12.&5.&10\end{array}\]}

From this, Articles 344 and 345 easily give the following distribution of all the roots \(\Omega\) into three periods of six terms, and of each of these into three periods of two terms: \[\Omega=(18,1)\left\{\begin{array}{l}(6,1)\left\{\begin{array}{ll}(2,1)\ldots&[1],[18]\\2,8)\ldots&[8],[11]\\(2,7)\ldots&[7],[12]\end{array}\right.\\[8pt](6,2)\left\{\begin{array}{ll}(2,2)\ldots&[2],[17]\\(2,16)\ldots&[3],[16]\\(2,14)\ldots&[5],[14]\end{array}\right.\\[8pt](6,4)\left\{\begin{array}{ll}(2,4)\ldots&[4],[15]\\(2,13)\ldots&[6],[13]\\(2,9)\ldots&[9],[10]\end{array}\right.\end{array}\right.\]

The equation \((A)\), whose roots are the sums \((6,1),\allowbreak (6,2),\allowbreak (6,4)\), is found to be \(x^3+xx-6x-7=0\), one of whose roots is found to be \(-1{,}2218761623\). Denoting this by \((6,1)\), we obtain
\SourceLeaf{440}
\[\begin{aligned}(6,2)&=4-(6,1)^2=2{,}5070186441\\(6,4)&=-5-(6,1)+(6,1)^2=-2.2851424818\end{aligned}\] Thus \(X\) will be resolved into three factors of degree 6 if these values are substituted in Article 348.

The equation \((B)\), whose roots are the sums \((2,1),\allowbreak (2.7),\allowbreak (2,8)\), is found to be \[x^3-(6,1)xx+((6,1)+(6,4))x-2-(6,2)=0\] or \[x^3+1{,}2218761623xx-3{,}5070186441x-4{,}5070186441=0\] One of its roots is found to be \(-1{,}3545631433\), which we will denote by \((2,1)\). The method of Article 346 gives the following equations, in which, for brevity, \(q\) is written for \((2,1)\): \[\begin{aligned}
(2,2)&=qq-2,\quad (2,3)=q^3-3q,\\
(2,4)&=q^4-4qq+2,\quad (2,5)=q^5-5q^3+5q\\
(2,6)&=q^6-6q^4+9qq-2,\\
(2,7)&=q^7-7q^5+14q^3-7q\\
(2,8)&=q^8-8q^6+20q^4-16qq+2,\\
(2,9)&=q^9-9q^7+27q^5-30q^3+9q
\end{aligned}\]

In the present case these equations can be developed more conveniently by the following considerations than by the rules of Article 346. If we suppose \[[1]=\cos\frac{kP}{19}+i\sin\frac{kP}{19}\] we have \[\begin{aligned}
\relax[18]&=\cos\frac{18kP}{19}+i\sin\frac{18kP}{19}\\
&=\cos\frac{kP}{19}-i\sin\frac{kP}{19},
\end{aligned}\] and therefore \((2,1)=2\cos\frac{kP}{19}\), and also, in general, \[[\lambda]=\cos\frac{\lambda kP}{19}+i\sin\frac{\lambda kP}{19},\] and therefore \((2,\lambda)=[\lambda]+[18\lambda]=[\lambda]+[-\lambda]=2\cos\frac{\lambda kP}{19}\) Thus, if \(\frac12q=\cos\omega\), we will have \((2,2)=2\cos2\omega,\allowbreak (2,3)=2\cos3\omega\), etc. The known equations for the cosines of multiple angles then give the same formulas as above. — These formulas now give the following numerical values: \[\begin{array}{r|r}(2,2)=-0{,}1651586909&(2,6)=0{,}4909709743\\(2,3)=1{,}5782810188&(2,7)=-1{,}7589475024\\(2,4)=-1{,}9727226068&(2,8)=1{,}8916344834\\(2,5)=1{,}0938963162&(2,9)=-0{,}8033908493\end{array}\]
\SourceLeaf{441}

The values of \((2,7),\allowbreak (2,8)\) can also be found from equation \((B)\), whose two remaining roots they are. The doubt as to which of these roots is \((2,7)\) and which is \((2,8)\) will be removed either by an approximate calculation using the preceding formulas or by sine tables. Even a cursory consultation of these shows that \((2,1)=2\cos\omega\) when we put \(\omega=\frac7{19}P\), whence we must have \[(2,7)=2\cos\frac{49}{19}P=2\cos\frac8{19}P,\] and \((2,8)=2\cos\frac{56}{19}P=2\cos\frac1{19}P\) Similarly, the sums \((2,2),\allowbreak (2,3),\allowbreak (2,5)\) can also be found by means of the equation \[x^3-(6,2)xx+((6,1)+(6,2))x-2-(6,4)=0\] whose roots they are; the uncertainty about which roots are to be set equal to the respective sums will be removed in exactly the same way as before. Likewise, the sums \((2,4),\allowbreak (2,6),\allowbreak (2,9)\) can also be found by the equation \[x^3-(6,4)xx+((6,2)+(6,4))x-2-(6,1)=0\]

Finally, \([1]\) and \([18]\) are the roots of the equation \(xx-(2,1)x+1=0\), one of which is \(=\frac12(2,1)+i\surd(1-\frac14(2,1)^2)=\frac12(2,1)+i\surd(\frac12-\frac14(2,2))\), the other \(=\frac12(2,1)-i\surd(\frac12-\frac14(2,2))\), giving the numerical values \(=-0{,}6772815716\pm0{,}7357239107i\). The sixteen remaining roots can be deduced either by developing the powers of either of these roots or by solving eight other similar equations. In the latter method, we must decide, either by sine tables or by the device to be explained in the next example, for which root the imaginary part is to take the positive sign and for which the negative sign. The following values were found in this way, with the upper sign understood to correspond to the first root and the lower to the second: \[\begin{aligned}\relax[1]\text{ and }[18]&=-0{,}6772815716\pm0{,}7357239107i\\\relax[2]\text{ and }[17]&=-0{,}0825793455\mp0{,}9965844930i\\\relax[3]\text{ and }[16]&=0{,}7891405094\pm0{,}6142127127i\\\relax[4]\text{ and }[15]&=-0{,}9863613034\pm0{,}1645945903i\\\relax[5]\text{ and }[14]&=0{,}5469481581\mp0{,}8371664783i\\\relax[6]\text{ and }[13]&=0{,}2454854871\pm0{,}9694002659i\\\relax[7]\text{ and }[12]&=-0{,}8794737512\mp0{,}4759473930i\\\relax[8]\text{ and }[11]&=0{,}9458172417\mp0{,}3246994692i\\\relax[9]\text{ and }[10]&=-0.4016954247\pm0{,}9157733267i\end{aligned}\]

\SourceLeaf{442}

\Needspace{6\baselineskip}\Article{354}
\textit{Second example, for \(n=17\).} Here we have \(n-1=2.2.2.2\), so the calculation of the roots \(\Omega\) is to be reduced to four quadratic equations. Here we will take the number 3 for the primitive root; its powers give the following least residues modulo 17: {\small\setlength{\arraycolsep}{2pt}\[\begin{array}{rrrrrrrrrrrrrrrr}0.&1.&2.&3.&4.&5.&6.&7.&8.&9.&10.&11.&12.&13.&14.&15\\1.&3.&9.&10.&13.&5.&15.&11.&16.&14.&8.&7.&4.&12.&2.&6\end{array}\]}

This gives the following distributions of the collection \(\Omega\) into two periods of eight terms, four of four terms, and eight of two terms: {\small\setlength{\arraycolsep}{2pt}\[\Omega=(16,1)\left\{\begin{array}{l}(8,1)\left\{\begin{array}{l}(4,1)\left\{\begin{array}{ll}(2,1)\ldots&[1],[16]\\(2,13)\ldots&[4],[13]\end{array}\right.\\[4pt](4,9)\left\{\begin{array}{ll}(2,9)\ldots&[8],[9]\\(2,15)\ldots&[2],[15]\end{array}\right.\end{array}\right.\\[8pt](8,3)\left\{\begin{array}{l}(4,3)\left\{\begin{array}{ll}(2,3)\ldots&[3],[14]\\(2,5)\ldots&[5],[12]\end{array}\right.\\[4pt](4.10)\left\{\begin{array}{ll}(2,10)\ldots&[7],[10]\\(2,11)\ldots&[6],[11]\end{array}\right.\end{array}\right.\end{array}\right.\]}

By the rules of Article 351, the equation \((A)\), whose roots are the sums \((8,1),\allowbreak (8,3)\), is found to be \(xx+x-4=0\). Its roots are calculated to be \(-\frac12+\frac12\surd 17=1{,}5615528128\) and \(-\frac12-\frac12\surd 17=-2{,}5615528128\). We will set the first \(=(8,1)\), so the second must necessarily be set \(=(8,3)\).

Further, the equation whose roots are the sums \((4,1)\) and \((4,9)\) is found to be \((B):xx-(8,1)x-1=0\). Its roots are \(\frac12(8,1)\pm\frac12\surd(4+(8,1)^2)=\frac12(8,1)\pm\frac12\surd(12+3(8,1)+4(8,3))\). We will set the root in which the radical quantity takes the positive sign, and whose numerical value is \(2{,}0494811777\), \(=(4,1)\); the other root, in which the radical quantity is taken negatively and whose value is \(-0{,}4879283649\), must then be denoted by \((4,9)\). The remaining sums of four terms, namely \((4,3)\) and \((4,10)\), can be found in two ways. \textit{First}, by the method of Article 346, which gives the following formulas, in which, for brevity, we write for \((4,1)\) the letter \(p\): \[\begin{aligned}(4,3)&=-\frac32+3p-\frac12p^3=0{,}3441507314\\(4,10)&=\frac32+2p-pp-\frac12p^3=-2{,}9057035442\end{aligned}\]
\SourceLeaf{443}

The same method also gives the formula \((4,9)=-1-6p+pp+p^3\), which yields the same value as we gave before. \textit{Second}, the sums \((4,3)\), \((4,10)\) can also be determined by solving the equation whose roots they are. This equation is \(xx-(8,3)x-1=0\), whence its roots are \(\frac12(8,3)\pm\frac12\surd(4+(8,3)^2)\), or \(\frac12(8,3)+\frac12\surd(12+4(8,1)+3(8,3))\) and \(\frac12(8,3)-\frac12\surd(12+4(8,1)+3(8,3))\). The doubt as to \textit{which} root should be denoted by \((4,3)\) and which by \((4,10)\) will be removed by the following device, which we mentioned in Article 352. Develop the product of \((4,1)-(4,9)\) and \((4,3)-(4,10)\); the result will be found to be \(2(8,1)-2(8,3)\)\HistoricalNote{*)}{The true nature of this device consists in the fact that it could be foreseen a priori that the developed product would contain no sums of four terms, but could be expressed through sums of eight terms alone. Those familiar with the subject will very easily see the reason for this, which must be passed over here for brevity.} Now the value of this expression is clearly positive, namely \(=+2\surd 17\), and the first factor of the product, \((4,1)-(4,9)\), is also positive, namely \(=+\surd(12+3(8,1)+4(8,3))\). Thus the other factor, \((4,3)-(4,10)\), must necessarily also be positive, and therefore \((4,3)\) must be set equal to the \textit{first} root, in which the radical is preceded by the positive sign, and \((4,10)\) to the second. This gives the same numerical values as above.

Having found all the sums of four terms, we proceed to the sums of two terms. The equation \((C)\), whose roots are \((2,1),\allowbreak (2,13)\), contained in \((4,1)\), is found to be \(xx-(4,1)x+(4,3)=0\). Its roots are \(\frac12(4,1)\pm\frac12\surd(-4(4,3)+(4,1)^2)\) or \(\frac12(4,1)\pm\frac12\surd(4+(4,9)-2(4,3))\). We set the one in which the radical quantity is taken positively, and whose value is found to be \(=1{,}8649444588\), \(=(2,1)\); thus \((2,13)\) will be equal to the other, whose value is \(=0{,}1845367189\). If we wish to find the remaining sums of two terms by the method of Article 346, for \((2,2),\allowbreak (2,3),\allowbreak (2,4),\allowbreak (2,5),\allowbreak (2,6),\allowbreak (2,7),\allowbreak (2,8)\) we can use the same formulas that we gave in the preceding example for quantities bearing the same notation, namely \((2,2)\), (or \((2,15)\)), \(=(2,1)^2-2\), etc. If, however, we prefer to calculate them in pairs by solving a quadratic equation, for \((2,9),\allowbreak (2,15)\) we find the equation \(xx-(4,9)x+(4,10)=0\), whose roots are \(\frac12(4,9)\pm\frac12\surd(4+(4,1)-2(4,10))\). How the choice of sign is to be determined here will be decided in the same way as above. Namely, developing the product of \((2,1)-(2,13)\) and \((2,9)-(2,15)\) gives \(-(4,1)+(4,9)-(4,3)+(4,10)\). Since this is clearly negative, whereas the factor \((2,1)-(2,13)\) is positive, \((2,9)-(2,15)\)
\SourceLeaf{444}
must necessarily be negative. Thus, in the expression given above, the upper, positive sign must be taken for \((2,15)\) and the lower, negative sign for \((2,9)\). From this we calculate \((2,9)=-1{,}9659461994\), \((2,15)=1{,}4780178344\). — Likewise, since developing the product of \((2,1)-(2,13)\) and \((2,3)-(2,5)\) gives \((4,9)-(4,10)\), hence a positive quantity, we conclude that the factor \((2,3)-(2,5)\) is positive. A calculation similar to the preceding one then gives \[\begin{aligned}(2,3)&=\tfrac12(4,3)+\tfrac12\surd(4+(4,10)-2(4,9))\\&=0{,}8914767116\\(2,5)&=\tfrac12(4,3)-\tfrac12\surd(4+(4,10)-2(4,9))\\&=-0{,}5473259801\end{aligned}\]

Finally, entirely similar operations give \[\begin{aligned}(2,10)&=\tfrac12(4,10)-\tfrac12\surd(4+(4,3)-2(4,1))\\&=-1{,}7004342715\\(2,11)&=\tfrac12(4,10)+\tfrac12\surd(4+(4,3)-2(4,1))\\&=-1{,}2052692728\end{aligned}\]

It remains to descend to the roots \(\Omega\) themselves. The equation \((D)\), whose roots are \([1]\) and \([16]\), is \(xx-(2,1)x+1=0\), whence the roots \(\frac12(2,1)\pm\frac12\surd((2,1)^2-4)\), or rather \(\frac12(2,1)\pm i\surd(4-(2,1)^2)\), or \(\frac12(2,1)\pm\frac12i\surd(2-(2,15))\). We take the upper sign for \([1]\), the lower for \([16]\). The fourteen remaining roots will be obtained either through the powers of \([1]\) or by solving seven quadratic equations, each of which gives a pair of roots; the uncertainty about the signs of the radical quantities can be removed by the same device as before. Thus \([4]\) and \([13]\) are the roots of the equation \(xx-(2,13)x+1=0\), and therefore \(\frac12(2,13)\pm\frac12i\surd(2-(2,9))\). Developing the product of \([1]-[16]\) and \([4]-[13]\), however, gives \((2,5)-(2,3)\), hence a negative real quantity. Since \([1]-[16]\) is \(+i\surd(2-(2,15))\), that is, the product of the imaginary quantity \(i\) and a \textit{positive} real quantity, \([4]-[13]\) must also be the product of \(i\) and a \textit{positive} real quantity, because \(ii=-1\). We conclude that the upper sign must be taken for \([4]\) and the lower for \([13]\). Similarly, for the roots \([8]\) and \([9]\) we find \(\frac12(2,9)\pm\frac12i\surd(2-(2,1))\). Since the product of \([1]-[16]\) and \([8]-[9]\) is \((2,9)-(2,10)\), hence negative, the upper sign must be taken for \([8]\) and the lower for \([9]\). Calculating the remaining roots in the same way, we obtain the following numerical values, where the upper signs are understood to correspond to the first roots and the lower to the second:
\SourceLeaf{445}
\[\begin{aligned}\relax[1],\ [16]&\quad\ldots\quad 0{,}9324722294\pm 0{,}3612416662i\\\relax [2],\ [15]&\quad\ldots\quad 0{,}7390089172\pm 0{,}6736956436i\\\relax [3],\ [14]&\quad\ldots\quad 0{,}4457383558\pm 0{,}8951632914i\\\relax [4],\ [13]&\quad\ldots\quad 0{,}0922683595\pm 0{,}9957341763i\\\relax [5],\ [12]&\quad\ldots\quad -0{,}2736629901\pm 0{,}9618256432i\\\relax [6],\ [11]&\quad\ldots\quad -0{,}6026346364\pm 0{,}7980172273i\\\relax [7],\ [10]&\quad\ldots\quad -0{,}8502171357\pm 0{,}5264321629i\\\relax [8],\ [9]&\quad\ldots\quad -0{,}9829730997\pm 0{,}1837495178i\end{aligned}\] \[\begin{matrix}*\qquad *\\{}*\end{matrix}\]

What has been given above could indeed suffice for solving the equation \mbox{\(x^n-1=0\)}, and therefore also for finding the trigonometric functions corresponding to arcs commensurable with the circumference. Yet, because of the importance of the subject, we cannot bring this investigation to an end without first adding here a few items from the wealth both of observations that shed light on it and of propositions related to it or dependent on it. We will chiefly choose those that can be completed without an extensive apparatus of other investigations, and we wish them to be regarded only as \textit{samples} of this very extensive theory, which is to be treated at length in the future.

